\documentclass[11pt,oneside]{article}

% ===============================
% Geometry and Layout
% ===============================
\usepackage[margin=1.15in]{geometry}
\usepackage{setspace}
\setstretch{1.15}

% ===============================
% Engine and Fonts (LuaLaTeX)
% ===============================
\usepackage{fontspec}
\usepackage{unicode-math}
\setmainfont{Libertinus Serif}
\setsansfont{Libertinus Sans}
\setmonofont{Libertinus Mono}
\setmathfont{Libertinus Math}

% ===============================
% Mathematics and Symbols
% ===============================
\usepackage{amsmath,amssymb}

% ===============================
% Microtypography and Links
% ===============================
\usepackage{microtype}
\usepackage{hyperref}

% ===============================
% Title
% ===============================
\title{Gallery Before Feed:\\
Why Engagement-Optimized Interfaces Destroy Coherence\\
and How Merge-Aware Systems Restore Meaning}
\author{Flyxion}
\date{\today}

\begin{document}
\maketitle

\begin{abstract}
Contemporary social media platforms present themselves as tools for expression, memory, and connection, yet their engagement-optimized architectures systematically undermine these functions. Users attempting to treat timelines as galleries, archives, or coherent identity surfaces encounter persistent interface intrusions, forced duplication, semantic fragmentation, and non-consensual content insertion. This paper argues that these failures are not incidental design flaws but necessary consequences of feed-based architectures that collapse semantic objects, events, and projections into a single irreversible stream.

We develop a gallery-first, merge-aware alternative grounded in event-sourced semantics, explicit object identity, and lossless consolidation. By separating persistent semantic objects from the events that reference them, treating feeds as projections rather than sources of truth, and formalizing merge as a first-class operation, the proposed architecture permits repetition without fragmentation and accumulation without drift. Interaction affordances are made latent rather than persistent, and interface states are explicitly separated into gallery, interaction, and annotation modes with formal transition rules.

The framework is situated in relation to conceptual blending, remix culture, and cognitive theories of analogy, showing how engagement-driven systems instantiate unconstrained blending as infrastructure and thereby guarantee semantic entropy. A formal appendix proves that semantic drift is inevitable under minimal feed-based assumptions, while semantic accumulation is guaranteed under merge-aware constraints.

The paper concludes that meaning cannot be optimized into existence through engagement metrics. It must be conserved by design. Gallery-first, merge-aware systems represent not a stylistic alternative to existing platforms, but a different architectural regime—one aligned with human sense-making, durable identity, and the preservation of meaning over time.
\end{abstract}

\section{Introduction}

Contemporary social media platforms increasingly present themselves as neutral tools for expression, memory, and social connection. Yet their interfaces betray a radically different purpose. What appears to the user as a personal timeline or gallery is, in practice, an instrument panel for engagement extraction, monetization, and behavioral reinforcement. The resulting tension is not aesthetic alone. It is structural.

This paper arises from a concrete and recurrent frustration: the attempt to use a personal social media timeline as a coherent visual archive or spaced repetition surface, only to be repeatedly disrupted by advertising affordances, forced duplication, algorithmic insertions, and non-removable content classes. The experience is not merely annoying. It is cognitively destructive.

The central claim of this work is that such frustration is not accidental, nor solvable through preference toggles or moderation tools. Instead, it reveals a deep incompatibility between engagement-optimized feed architectures and the requirements of coherent personal memory systems. In response, this paper proposes an alternative paradigm: a gallery-first, merge-aware, event-sourced social substrate in which meaning accretes rather than fragments.

\section{The Timeline as Monetization Surface}

A defining feature of contemporary platforms is the presence of persistent monetization controls embedded directly into user-created content. Buttons labeled with imperatives such as “Boost Post” or “See Ads and Insights” are rendered with visual priority, often in primary colors that dominate the surrounding composition. These elements are not optional overlays. They are inseparable from the post itself.

This design choice reclassifies personal expression as latent inventory. A photograph, rather than being treated as an artifact to be viewed, remembered, or reflected upon, is framed as an underperforming asset awaiting optimization. The interface does not ask whether the user wishes to sell or promote. It assumes that monetization is the natural continuation of expression.

From a systems perspective, this introduces a category error. A gallery is an archive of artifacts. An advertising dashboard is a control surface for conversion optimization. By collapsing these into a single interface layer, the platform destroys the semantic integrity of both. The user cannot view their own work without being reminded that it has failed to perform as a revenue-generating unit.

The aesthetic damage is immediate, but the deeper harm lies in the erosion of authorship. The creator is subtly repositioned as a marketer of their own life, regardless of intent.

\section{Feed-Ordered Atoms and the Incoherence of Duplication}

Another structural failure emerges when identical or near-identical content is posted more than once. A user may post the same image with a minor textual variation, such as the addition of a hashtag, or may unintentionally create multiple timeline entries referencing the same artifact. In such cases, the platform offers no reconciliation mechanism.

Each post is treated as an immutable, feed-ordered atom. Reactions, comments, and engagement metrics attach to the publication event rather than to the underlying content. The system lacks any notion of semantic equivalence, versioning, or mergeability. As a result, identical images compete with themselves for attention and validation, fragmenting both meaning and response.

The user intuitively recognizes that these entries refer to the same object and should be unified. The platform, however, recognizes only chronological difference. This reveals an architectural assumption: that novelty and temporal ordering matter more than coherence or identity.

The inability to merge duplicate posts is not an oversight. It reflects a deeper refusal to acknowledge that meaning persists beyond momentary exposure.

\section{Identity Surfaces and the Cover Photo Paradox}

The problem becomes even more pronounced when an image occupies multiple identity surfaces simultaneously, such as appearing both as a cover photo and as one or more timeline posts. These representations are treated as unrelated entities. They maintain separate engagement counts, separate contexts, and separate lifecycles.

There is no mechanism by which the system can recognize that these surfaces reference the same semantic object. The result is identity fragmentation. The same image exists in parallel realities, each accumulating its own partial history, none of which fully represents the whole.

For a user attempting to maintain a coherent visual identity or archive, this fragmentation produces a sense of disorder and loss of control. The platform forces multiplicity without offering synthesis. Identity becomes a scatterplot rather than a trajectory.

\section{Algorithmic Intrusions and Non-Removable Content Classes}

Perhaps the most psychologically disruptive feature of modern platforms is the forced insertion of algorithmically selected content classes, such as short-form video reels, into spaces that users experience as personal. These elements persist even when the user does not engage with them, does not follow their creators, and repeatedly signals disinterest.

The reason for this persistence is structural. Such content classes are not personalized recommendations in the traditional sense. They are strategic instruments designed to recalibrate engagement baselines, compete with external platforms, and exploit high-stability attention triggers. Their presence is not negotiable because they are not subordinate to user intent.

Blocking individual creators or hiding specific items fails to converge because the system optimizes at the level of content archetypes rather than sources. The user is not being shown something by mistake. They are being exposed to an equilibrium state.

This breaks the illusion of agency. The platform no longer responds meaningfully to preference signals. Instead, it enforces exposure.

\section{Why These Failures Are Systemic}

Taken together, these phenomena reveal a consistent pattern. The platform does not treat posts as durable artifacts, users as authors, or timelines as archives. It treats everything as transient input to an optimization loop.

Engagement metrics replace memory. Chronology replaces identity. Exposure replaces intention. Under such conditions, any attempt to use the system as a gallery, notebook, or spaced repetition surface will fail.

The frustration experienced by users who desire coherence is not emotional excess. It is a rational response to a system that systematically violates the principles of persistence, composability, and semantic continuity.

\section{Gallery-First Systems}

A gallery-first system inverts these priorities. Instead of treating posts as feed entries, it treats them as references to persistent semantic objects. Images, texts, and other artifacts possess stable identities independent of when or how often they are referenced.

In such a system, the primary interface is not a feed but a navigable archive. Chronology becomes a secondary projection rather than the organizing principle. Artifacts can be revisited, revised, and recontextualized without duplication.

Importantly, visibility and repetition strengthen an object’s history rather than fragmenting it. Re-posting an image does not create a new object. It creates a new event attached to an existing one.

\section{Merge-Aware Semantics}

Merge awareness is the logical extension of this model. When two posts reference the same artifact, the system recognizes their equivalence and offers unification. Reactions, comments, and annotations are aggregated into a shared history.

Merging does not erase provenance. Instead, it preserves the event graph while eliminating semantic redundancy. The user retains the ability to see when and how an artifact was resurfaced, while benefiting from a consolidated representation.

This directly supports spaced repetition. Returning to an artifact reinforces it rather than duplicating it. Memory becomes cumulative rather than noisy.

\section{Event-Sourced Identity}

Under this architecture, identity is not a collection of posts but a trajectory through an event space. Each interaction leaves a trace. Nothing is overwritten, yet nothing is forced to fragment.

This model aligns with event-sourced systems in computing, where state is derived from history rather than stored directly. It also aligns with human cognition, which builds meaning through revisitation, consolidation, and reinterpretation.

\section{Conclusion}

The failures of contemporary social media interfaces are not matters of taste or preference. They are the predictable consequences of optimizing for engagement without regard for coherence. When platforms treat personal expression as ad inventory and memory as a feed, they inevitably destroy both.

A gallery-first, merge-aware system restores what engagement optimization erases. It treats artifacts as durable, identity as historical, and repetition as reinforcement rather than noise. Such a system does not forbid monetization or discovery, but it subordinates them to meaning rather than allowing them to dominate.

The frustration that motivates this paper is therefore diagnostic. It signals a mismatch between human sense-making and machine optimization. The solution is not better filters or stronger moderation, but a different foundation entirely.

Meaning cannot be optimized into existence. It must be conserved by design.

\section{Formal Mapping to Merge-Aware Event Primitives}

This section makes explicit the correspondence between the proposed gallery-first system and a concrete set of formal primitives. The goal is not to introduce an implementation, but to define a minimal semantic substrate capable of supporting coherence, mergeability, and identity persistence. The notation is intentionally conservative and compatible with event-sourced and proof-carrying systems.

\subsection{Semantic Objects}

Let $\mathcal{O}$ denote the set of semantic objects. A semantic object $o \in \mathcal{O}$ represents a persistent artifact such as an image, text, or composite media object. Each object possesses a stable identity independent of presentation context.

Formally, an object is defined as
\[
o := \langle \mathrm{id}_o, \mathrm{payload}_o \rangle
\]
where $\mathrm{id}_o$ is a globally unique identifier and $\mathrm{payload}_o$ is immutable content. Mutations to content are represented not by overwriting $\mathrm{payload}_o$, but by the creation of successor objects with explicit lineage.

\subsection{Events}

Let $\mathcal{E}$ denote the set of events. An event represents an irreversible action taken with respect to one or more semantic objects. Events do not store state; they produce history.

An event $e \in \mathcal{E}$ is defined as
\[
e := \langle \mathrm{id}_e, \mathrm{type}_e, \mathrm{targets}_e, t_e \rangle
\]
where $\mathrm{id}_e$ is a unique event identifier, $\mathrm{type}_e$ specifies the event class, $\mathrm{targets}_e \subseteq \mathcal{O}$ is the set of objects referenced by the event, and $t_e$ is a timestamp.

Crucially, multiple events may target the same object. Reposting, resurfacing, annotating, or recontextualizing an image all produce distinct events referencing a shared semantic object rather than duplicating it.

\subsection{Timelines as Projections}

A timeline is not a primary data structure but a projection over the event graph. Let $\mathcal{G} = (\mathcal{O}, \mathcal{E})$ denote the bipartite object–event graph. A timeline $\tau$ is a view defined by a projection operator
\[
\tau = \pi(\mathcal{G}, \theta)
\]
where $\theta$ specifies ordering, filtering, or contextual constraints.

Under this model, deleting a timeline entry does not delete an object or erase history; it merely alters a projection. Identity coherence is therefore preserved independently of presentation.

\subsection{Merge Operations}

Merge awareness arises from the ability to recognize equivalence between events that reference the same semantic object. Let $e_1, e_2 \in \mathcal{E}$ such that
\[
\mathrm{targets}_{e_1} = \mathrm{targets}_{e_2} = \{o\}
\]
for some object $o$.

A merge operation $\mu$ produces a consolidated event view without destroying provenance:
\[
\mu(e_1, e_2) := \langle \mathrm{id}_\mu, \{\mathrm{id}_{e_1}, \mathrm{id}_{e_2}\}, o \rangle
\]

Reactions, annotations, and metadata associated with $e_1$ and $e_2$ are lifted to the merged structure while retaining references to their originating events. The merge is therefore associative and idempotent, but not destructive.

\subsection{Reaction Aggregation}

Let $\mathcal{R}$ denote the set of reactions. In feed-based systems, reactions bind to posts. In a merge-aware system, reactions bind to objects via events.

A reaction is defined as
\[
r := \langle \mathrm{actor}, o, e, t_r \rangle
\]
indicating that an actor reacted to object $o$ through event $e$ at time $t_r$.

Aggregated reaction counts are derived quantities:
\[
\mathrm{count}(o) = |\{ r \in \mathcal{R} \mid r.o = o \}|
\]
Thus, reposting an object strengthens its reaction history rather than splitting it across duplicates.

\subsection{Cover Surfaces as Identity Anchors}

Persistent identity surfaces such as cover images are modeled as bindings rather than copies. A cover surface is represented by a binding relation
\[
\beta : \mathrm{Surface} \rightarrow \mathcal{O}
\]
which associates a semantic object with an identity role.

Changing a cover image replaces the binding, not the object. The object’s event history remains intact and accessible from all projections.

\subsection{Suppression of Non-Consensual Insertions}

Algorithmic insertions such as reels are treated as events without explicit user consent. In a gallery-first system, all events must reference an object explicitly accepted into the user’s object space.

Formally, let $\mathcal{O}_u \subseteq \mathcal{O}$ denote the object set authorized by user $u$. An event $e$ is admissible in projections for $u$ only if
\[
\mathrm{targets}_e \subseteq \mathcal{O}_u
\]

This constraint eliminates forced content classes by construction rather than moderation.

\subsection{Interpretation}

This mapping demonstrates that gallery-first behavior is not an aesthetic preference but the consequence of a different primitive set. By separating objects from events, treating timelines as projections, and making merge operations first-class, the system restores coherence, reversibility of views, and accumulation of meaning.

Feed-based incoherence is not emergent complexity. It is the direct result of collapsing objects, events, and projections into a single irreversible structure. The primitives defined here explicitly refuse that collapse.

\section{Relation to Conceptual Blending}

The gallery-first, merge-aware model bears an immediate resemblance to the theory of conceptual blending developed in cognitive science, most notably by \textbf{0} and \textbf{1}. Both frameworks reject purely symbolic substitution and instead emphasize the construction of meaning through structured combination. However, the similarity is partial and instructive, because the proposed system can be understood as a constrained, infrastructural specialization of blending rather than a general cognitive model.

\subsection{Conceptual Blending in Brief}

Conceptual blending theory describes how humans integrate multiple mental spaces into a blended space that exhibits emergent structure. Inputs retain partial identity, shared structure is projected into a generic space, and the blend produces novel inferences not present in any single input. Crucially, blends are not simple unions. They selectively compress, align, and elaborate structure.

Blending is therefore creative, opportunistic, and context-sensitive. It operates without a fixed notion of identity persistence, auditability, or reversibility. Once a blend is formed, its inputs may no longer be recoverable in a precise or formal way.

\subsection{Semantic Objects as Stabilized Mental Spaces}

In the gallery-first system, semantic objects play a role analogous to mental spaces. Each object encapsulates a coherent bundle of meaning, content, and identity. Unlike mental spaces, however, semantic objects are intentionally stabilized. They possess immutable identifiers and explicit boundaries.

This stabilization is a deliberate departure from cognitive blending. Whereas blending thrives on fluidity, the gallery-first system enforces persistence as an infrastructural constraint. The purpose is not to model cognition directly, but to support memory, authorship, and accountability in a shared computational environment.

\subsection{Merge as Constrained Blending}

The merge operation resembles blending in that it allows multiple event contexts to be unified into a higher-level structure. Two reposts of the same image, for example, can be merged into a consolidated representation that aggregates reactions and annotations. This resembles the compression and integration characteristic of blends.

The key difference lies in constraint. A merge does not invent new semantic content. It does not produce emergent meaning beyond what is already present in the referenced object and its event history. Instead, it performs a lossless consolidation. Provenance is preserved, reversibility is maintained, and no implicit inference is introduced.

In this sense, merge can be seen as a restricted form of blending in which creativity is intentionally suppressed in favor of coherence.

\subsection{Event Graphs versus Emergent Blends}

Conceptual blending allows the same inputs to yield different blends depending on context, goals, and attentional framing. The gallery-first model explicitly resists this multiplicity. Events referencing the same object do not generate divergent semantic offspring. They accumulate into a single event graph.

This design choice reflects a different priority. Conceptual blending optimizes for sense-making under uncertainty. Gallery-first systems optimize for durable meaning under repetition. Where blending embraces ambiguity, merge-aware systems bound it.

\subsection{Why Unconstrained Blending Fails as Infrastructure}

If conceptual blending were used directly as an infrastructural principle for social systems, the result would be catastrophic for identity coherence. Blends do not guarantee traceability. They do not enforce equivalence relations. They do not provide canonical representations.

This is precisely the failure mode observed in engagement-driven platforms. Content is continually recontextualized, recombined, and abstracted without stable anchors. The feed becomes a perpetual blend space in which nothing settles long enough to accumulate meaning.

The gallery-first model can therefore be understood as a corrective response: it retains the insight that meaning emerges from combination, while refusing to allow combination to dissolve identity.

\subsection{Interpretive Summary}

Conceptual blending explains how humans generate new meaning. Gallery-first, merge-aware systems explain how meaning can be preserved once generated. The former is generative and cognitive. The latter is conservative and infrastructural.

The relationship between the two is complementary rather than competitive. Conceptual blending operates upstream, in the formation of ideas and interpretations. Merge-aware galleries operate downstream, in the maintenance of artifacts, histories, and shared reference points.

By constraining blending into explicit merge operations over persistent objects, the system reconciles creativity with coherence. It allows repetition without fragmentation, integration without erasure, and growth without semantic drift.

\section{Unbounded Blending, Remix Culture, and the Inevitable Drift of Feed-Based Systems}

The failure modes diagnosed in engagement-optimized platforms can be understood more precisely by situating them between two related but distinct conceptual frameworks: conceptual blending in cognitive science and remix culture in networked media. Both illuminate why feed-based systems drift toward incoherence, and both clarify why a gallery-first, merge-aware architecture represents a fundamentally different regime rather than a refinement of existing designs.

\subsection{Conceptual Blending as an Unconstrained Process}

Conceptual blending, as articulated by \textbf{0} and \textbf{1}, describes how meaning emerges through the integration of multiple mental spaces into blended spaces that exhibit novel structure. This process is opportunistic, selective, and context-dependent. Inputs are partially projected, compressed, and elaborated, producing inferences not present in the originals.

Crucially, conceptual blending does not preserve identity in a strict sense. Mental spaces are temporary, task-relative, and often unrecoverable once a blend stabilizes. This is a feature, not a bug, at the level of cognition. It allows humans to generate insight, metaphor, and creativity under uncertainty.

However, when this same unconstrained blending dynamic is instantiated as infrastructure rather than cognition, it becomes pathological. Without stable anchors, equivalence relations, or canonical forms, blended structures cannot accumulate durable meaning. They continuously overwrite and reframe one another.

\subsection{Feeds as Persistent Blend Spaces}

Engagement-driven feeds effectively function as perpetual blend spaces. Each item is stripped of stable identity and reintroduced as a momentary stimulus, recombined with adjacent content through scrolling context, algorithmic juxtaposition, and visual similarity.

In such systems, repetition does not reinforce meaning. It generates divergence. The same image, phrase, or theme appears in multiple blended contexts, each time acquiring slightly different associations, reactions, and interpretations. Because there is no underlying semantic object to which these variations can attach, the system cannot converge.

This explains why attempts to “fix” feeds through better recommendation tuning fail. The problem is not the quality of blends but the absence of constraints that would allow blends to settle into shared reference.

\subsection{Remix Culture as a Cultural Parallel}

Remix culture provides a cultural analogue to this dynamic. In remix ecosystems, artifacts are copied, modified, recontextualized, and redistributed with minimal friction. Attribution may be partial or lost entirely. Lineage becomes aesthetic rather than structural.

This mode of production is generative and expressive, but it relies on an implicit substrate of abundance and disposability. Artifacts are expected to mutate and proliferate. Persistence is optional. Canonical versions are often rejected as authoritarian or creatively stifling.

Feed-based platforms implicitly adopt remix culture as an infrastructural default. Every post is treated as remixable, transient, and replaceable. Even original content is immediately absorbed into an endless chain of reinterpretation and abstraction.

\subsection{Why Remix Logic Fails as a Memory System}

While remix culture excels at cultural production, it fails catastrophically as a memory system. Without mechanisms for consolidation, equivalence, and merge, meaning diffuses faster than it accumulates. Recognition replaces understanding. Virality replaces recall.

This is why users attempting to use feeds as galleries, notebooks, or personal archives experience increasing frustration. They are attempting to perform conservation in a system optimized for dispersion.

The gallery-first model explicitly rejects remix logic as a foundational principle. It does not forbid remix, but it requires remix to be represented as an event referencing a stable object. Transformation becomes explicit rather than implicit. Lineage becomes structural rather than stylistic.

\subsection{Merge-Aware Systems as Bounded Blending}

Merge-aware galleries can be understood as implementing bounded blending. Combination is permitted, but only under constraints that preserve identity, provenance, and reversibility. Events accumulate rather than overwrite. Repetition strengthens rather than fragments.

Where unconstrained blending generates novelty, bounded blending generates coherence. Where remix culture privileges proliferation, merge-aware systems privilege consolidation.

This distinction is not aesthetic. It is thermodynamic. Unbounded blending increases semantic entropy. Merge-aware accumulation reduces it by channeling variation into structured history.

\subsection{Theorem-Style Claim: Drift Is Inevitable in Feed Architectures}

The preceding analysis supports a strong claim.

Any system in which content units lack stable identity, reactions bind to publication events rather than semantic objects, and combination occurs implicitly through context rather than explicitly through merge operations will necessarily exhibit unbounded semantic drift.

This drift is not a failure of moderation, recommendation, or governance. It is a structural consequence of collapsing objects, events, and projections into a single feed-ordered primitive.

Conversely, any system that separates semantic objects from events, treats feeds as projections rather than sources of truth, and provides explicit merge operations will necessarily permit the accumulation of meaning over time.

\subsection{Synthesis}

Conceptual blending explains how meaning is generated. Remix culture explains how meaning propagates. Feed-based platforms confuse these processes with preservation and therefore fail at it.

Gallery-first, merge-aware systems occupy a different niche. They are not engines of novelty, nor are they neutral conduits of remix. They are conservation systems for meaning in environments where repetition is inevitable.

By bounding blending and formalizing merge, such systems allow creativity upstream and coherence downstream. They do not eliminate drift by suppressing expression. They eliminate it by refusing to mistake motion for meaning.

\section{Surfaces, Essences, and the Failure of Instance-Bound Meaning}

The distinction between surface presentation and underlying meaning has been articulated most clearly in the cognitive work of \textbf{0} and Emmanuel Sander, particularly in their analysis of analogy as the primary mechanism of human thought. Their central claim is that meaning does not reside in surface forms themselves, but in the recognition that multiple surface encounters refer to the same underlying conceptual essence.

In this view, repetition is not redundancy. Encountering different surface expressions of the same idea strengthens understanding precisely because the mind recognizes their shared referent. Meaning accumulates when the system can reliably map surface variation onto a stable conceptual core.

This distinction maps cleanly onto the architectural separation between semantic objects and events developed in this paper. Semantic objects correspond to essences: persistent, identity-bearing referents. Events and projections correspond to surfaces: contextual, temporal, and variable presentations of those referents.

Feed-based systems collapse this distinction. Each post is treated as a new surface without an acknowledged essence beneath it. Even when two posts contain identical payloads, the system fails to recognize their equivalence. As a result, repetition does not reinforce meaning. It fractures it. The user experiences the same image or idea as multiple unrelated instances rather than as multiple views of a single object.

From the perspective of Hofstadter and Sander’s framework, this represents a categorical error. The system prevents the formation of analogy by refusing to acknowledge sameness across contexts. Without a stable essence, the mind cannot perform the conceptual compression that gives rise to meaning. Surface variation becomes noise rather than signal.

Merge-aware systems correct this error by making essences explicit. When multiple events reference the same semantic object, the system affirms that these surfaces belong together. The user is not required to infer equivalence manually. The infrastructure performs the analogical alignment on their behalf.

This is not a cosmetic improvement. It changes the cognitive contract between user and system. Instead of forcing the user to continually reconstruct meaning from fragmented surfaces, the system supports recognition by preserving identity across repetition.

Seen this way, merge operations are not merely technical conveniences. They are formal analogical commitments. They assert that two encounters are “about the same thing,” and they preserve that assertion over time. Reaction aggregation, event consolidation, and gallery-first projections all follow naturally from this commitment.

The persistent intrusion of prompts, metrics, and engagement affordances in feed-based interfaces further exacerbates the problem by competing with essence recognition. Attention is drawn to performance signals rather than to the object itself. The surface overwhelms the referent.

A gallery-first system, by contrast, minimizes surface noise in its default state. By suppressing interaction affordances until explicitly invoked, it allows the user to encounter the artifact as an object rather than as an instrument. This creates the conditions under which essence recognition can occur.

In summary, the failure of feed-based platforms is not merely that they optimize engagement. It is that they systematically undermine the cognitive mechanism by which meaning stabilizes. By binding reactions to surfaces rather than to essences, they prevent analogy from doing its work.

A merge-aware, gallery-first architecture can therefore be understood as analogy-respecting infrastructure. It aligns system behavior with the way humans naturally form and maintain meaning, not by simulating cognition, but by refusing to sabotage it.

\section{Design Corollary: Interface Elements as Architectural Commitments}

The preceding sections argue that coherence or incoherence is not an emergent property of content, but a direct consequence of infrastructural primitives. This claim has an immediate design corollary: interface elements are not cosmetic. They are commitments. Every visible affordance reveals what the system believes a post is for.

A useful illustration is found in the default timeline composer and gallery views of contemporary platforms. In what is nominally presented as a personal timeline or archive, the interface inserts a persistent, stylized prompt reading “Write a note,” often rendered as a speech bubble or callout. This element is visually dominant, stylistically incongruent with photographic content, and permanently present regardless of user intent.

From a design perspective, this element encodes a strong assumption: that the primary function of the timeline is conversational emission rather than archival presentation. The timeline is treated as a prompt surface, not a gallery surface. Silence, curation, and restraint are not first-class states.

For users attempting to maintain a portfolio-like or gallery-first presentation, this single affordance is sufficient to break coherence. The presence of an unsolicited input prompt transforms the page from an exhibition into a workshop that never closes. The system does not permit the user to say, structurally, “I am done speaking here.”

This problem generalizes. Like and comment buttons persist beneath every image, regardless of whether interaction is desired or appropriate. These controls visually anchor attention away from the artifact itself and toward its engagement performance. The image is never allowed to simply be seen. It must also be evaluated.

A revealing counterexample already exists within the same platform. When posts are archived, the interface often suppresses visible engagement controls. Images appear without reaction counters or comment affordances. The result is immediately perceptible: the presentation becomes calmer, more legible, and closer to a true gallery. The content has not changed. Only the projection has.

This demonstrates a crucial point. The system already possesses the capacity to separate artifact viewing from interaction mechanics. The failure is not technical but philosophical. Engagement is assumed to be the default mode of perception rather than an optional action.

The gallery-first corollary is therefore simple but radical: interaction affordances should be latent, not persistent. Like, comment, promote, or annotate controls should appear only after an explicit interaction gesture, such as tapping, focusing, or entering a deliberate engagement mode. In the absence of such a gesture, the artifact should be presented without commentary infrastructure.

Formally, this corresponds to treating interaction as an event layered atop an object, not as a permanent property of its display. The default projection of an object should be contemplative, not transactional. Engagement should be possible everywhere but visible nowhere until invoked.

Under this model, the “Write a note” prompt is not merely misplaced; it is conceptually invalid. Creation and curation are distinct modes. A system that collapses them forces users into perpetual emission, undermining the possibility of reflection, sequencing, or completion.

The design corollary extends further. If merge-aware semantics are implemented, then repeated postings of the same object strengthen its presence without multiplying interface clutter. If interaction controls are deferred, then galleries remain visually stable even as engagement accumulates. If timelines are treated as projections rather than primary storage, then users may choose views optimized for memory, presentation, or conversation without corrupting the underlying structure.

What appears, then, as an aesthetic complaint about an ugly speech bubble is in fact a precise diagnosis. The interface reveals that the system does not recognize a gallery state as legitimate. It does not allow an artifact to rest. Everything must ask for more.

A gallery-first system makes the opposite commitment. It assumes that meaning often emerges through restraint. It allows objects to exist without solicitation. It treats interaction as intentional rather than compulsory.

This is not a matter of taste. It is the visible surface of an architectural choice about whether a system optimizes for engagement or for coherence. Only one of these permits memory to form.

\section{UI State Taxonomy: Gallery, Interaction, and Annotation Modes}

A gallery-first, merge-aware system requires an explicit separation of interface states. The absence of such separation is one of the primary reasons engagement-optimized platforms collapse presentation, interaction, and solicitation into a single incoherent surface.

Three states are sufficient to restore coherence.

The \emph{gallery state} is the default. In this state, semantic objects are presented without visible interaction affordances. Images appear without like counts, comment prompts, share buttons, or monetization controls. The interface assumes contemplation rather than response. Navigation is spatial or archival rather than conversational. This state corresponds to a projection optimized for recognition, memory, and aesthetic continuity.

The \emph{interaction state} is entered deliberately. It may be triggered by a tap, click, focus gesture, or explicit mode switch. Only in this state do interaction affordances become visible. Likes, comments, reactions, and annotations appear as overlays or contextual panels. Importantly, their appearance signals a change in intent: the user has chosen to act rather than to view.

The \emph{annotation state} is distinct from interaction. Annotation is not reaction. It is the act of adding structure, explanation, or extension to an object’s history. Entering annotation mode creates a new event explicitly tied to the object, not a transient conversational reply. Annotations are therefore durable, merge-aware, and discoverable across projections.

Transitions between these states must be explicit and reversible. The system should never assume that viewing implies responding, or that responding implies editing. By separating these modes, the interface allows the same semantic object to serve as an artwork, a discussion locus, or a working document without duplicating or degrading it.

Feed-based systems lack this taxonomy entirely. Every surface is simultaneously gallery, comment thread, advertisement, and prompt. The result is perpetual mode confusion. Gallery-first systems resolve this by treating state as a first-class concept rather than an emergent side effect.

\section{The Fear of Silence: Empty States and Engagement Anxiety}

Modern platforms exhibit a pervasive discomfort with silence. Empty states are treated not as legitimate resting points, but as failures requiring immediate remediation. This anxiety manifests in prompts, nudges, and placeholders that insist the user speak, react, or scroll.

The “Write a note” prompt embedded in personal timelines is a canonical example. It occupies visual space even when the user has nothing to say. Its presence asserts that a lack of output is an error condition. The interface does not permit quiet curation. It demands emission.

This design pattern reveals an underlying assumption: value is generated only through activity. A silent gallery is interpreted as an underperforming feed. From an engagement perspective, this makes sense. From a memory or portfolio perspective, it is catastrophic.

Silence is not absence. In archival systems, silence is completion. A finished gallery does not ask questions of the viewer. It presents itself and waits. By contrast, engagement-driven interfaces cannot wait. Waiting does not produce metrics.

The suppression of engagement controls in archived views inadvertently demonstrates this truth. When likes and comments disappear, the content breathes. The user experiences relief not because functionality was lost, but because solicitation was removed. The system briefly permits the artifact to exist without demanding a response.

Empty states, in a gallery-first system, are therefore not gaps to be filled but signals that a projection is complete. They communicate sufficiency. They tell the user: nothing is required of you here.

Engagement platforms resist this because silence is not legible to optimization loops. It produces no signal. But meaning often does not produce signal either, at least not immediately. It accumulates slowly, through revisitation and recognition.

A system that cannot tolerate silence cannot support memory. It will always privilege noise over coherence. By contrast, a system that treats silence as a valid state creates room for reflection, sequencing, and eventual return.

The design implication is straightforward. Prompts should be optional, latent, or explicitly summoned. Default surfaces should be allowed to remain visually and semantically quiet. The absence of interaction should not be framed as a failure.

Silence is not disengagement. It is often the precondition for meaning.

\section{Formal Transition Rules Between Interface States}

To ensure that gallery-first, merge-aware behavior is not merely aspirational, the system must define explicit and enforceable transition rules between interface states. These rules determine when interaction affordances appear, how events are generated, and how projections change without corrupting underlying semantic structure.

Let $\mathcal{S} = \{G, I, A\}$ denote the set of interface states, where $G$ is the gallery state, $I$ the interaction state, and $A$ the annotation state. Let $\mathcal{O}$ denote the set of semantic objects and $\mathcal{E}$ the set of events as previously defined.

\subsection{State Invariants}

Each state enforces invariants that constrain visible affordances and admissible actions.

In the gallery state $G$, semantic objects are rendered without interaction controls. No new events may be generated except navigation events, which do not target objects. The gallery state is observational and non-productive by design.

In the interaction state $I$, reaction and conversational affordances are visible. Events generated in this state must reference an existing semantic object and must not modify object identity or payload.

In the annotation state $A$, users may generate durable semantic extensions. Events generated in this state explicitly alter the object–event graph by adding structured annotations, links, or revisions.

These invariants prevent accidental escalation. Viewing does not imply reacting, reacting does not imply annotating, and annotating does not imply editing original content.

\subsection{Transition Function}

Define a transition function
\[
\delta : \mathcal{S} \times \mathcal{U} \rightarrow \mathcal{S}
\]
where $\mathcal{U}$ is the set of admissible user actions.

Transitions are governed by intent-bearing actions rather than passive gestures. For example, scrolling or passive focus does not alter state.

Valid transitions include
\[
\delta(G, \text{select}(o)) = I
\]
\[
\delta(I, \text{annotate}) = A
\]
\[
\delta(I, \text{dismiss}) = G
\]
\[
\delta(A, \text{commit}) = I
\]
\[
\delta(A, \text{cancel}) = I
\]

There is no direct transition from $G$ to $A$. Annotation requires an explicit intermediate interaction state, ensuring that durable changes are never initiated accidentally.

\subsection{Event Generation Rules}

Let $e \in \mathcal{E}$ be an event generated by user action $u$ in state $s$.

An event is admissible if and only if
\[
\mathrm{state}(u) = s \;\wedge\; \mathrm{type}(e) \in \mathrm{AllowedEvents}(s)
\]

In particular,
\[
\mathrm{AllowedEvents}(G) = \varnothing
\]
\[
\mathrm{AllowedEvents}(I) = \{\text{react}, \text{comment}, \text{share}\}
\]
\[
\mathrm{AllowedEvents}(A) = \{\text{annotate}, \text{link}, \text{revise}\}
\]

This formalizes a key design commitment: the default experience produces no semantic noise. All meaning-bearing actions are intentional.

\subsection{Projection Consistency}

Let $\pi_s$ denote the projection associated with state $s$. For any object $o \in \mathcal{O}$ and any two states $s_1, s_2 \in \mathcal{S}$, the following must hold:
\[
\mathrm{id}_o(\pi_{s_1}) = \mathrm{id}_o(\pi_{s_2})
\]

State transitions alter visibility and affordances but never object identity or history. This guarantees that changing modes does not fracture meaning or duplicate content.

\subsection{Merge Preservation Under Transition}

Merge operations are orthogonal to interface state. If events $e_1, e_2$ are mergeable under the criteria defined previously, then for all states $s \in \mathcal{S}$,
\[
\mu(e_1, e_2) \text{ is invariant under } \pi_s
\]

Thus, merging is a semantic operation, not a presentational one. Interface state may affect how merged history is displayed, but not whether it exists.

\subsection{Interpretation}

These transition rules enforce a strict separation between perception and action. They ensure that silence is a valid and stable state, that interaction is deliberate, and that semantic modification is explicit.

In feed-based systems, transitions are implicit and continuous. Every scroll is an interaction. Every view is a potential signal. By contrast, gallery-first systems treat state transitions as commitments. Nothing happens unless the user chooses for it to happen.

This formalization makes accidental engagement impossible by construction. It encodes restraint as a system property rather than a behavioral expectation.

\section{Conclusion: Coherence as a First-Class System Property}

This paper began with an everyday frustration: the inability to treat a personal timeline as a coherent gallery without constant interference from prompts, metrics, and unsolicited content. That frustration proved diagnostic. When examined carefully, it revealed not a failure of taste or configuration, but a deep architectural incompatibility between engagement-optimized systems and the requirements of durable meaning.

Through successive layers of analysis, this work has shown that feed-based platforms collapse semantic objects, events, projections, and interaction states into a single undifferentiated surface. This collapse produces unavoidable consequences. Identity fragments. Repetition generates noise rather than reinforcement. Interaction becomes compulsory rather than intentional. Silence is treated as error. Drift is not merely possible but inevitable.

The proposed gallery-first, merge-aware system demonstrates that these outcomes are not inherent to digital media. They follow from specific design commitments. By separating semantic objects from events, treating timelines as projections rather than sources of truth, and making merge operations explicit and lossless, the system restores accumulation where feeds enforce churn.

The comparison with conceptual blending clarified an important distinction. Human cognition relies on unconstrained blending to generate insight, metaphor, and novelty. Infrastructure, however, cannot rely on unconstrained blending without dissolving identity. Feed architectures mistakenly instantiate a cognitive process as a storage substrate, guaranteeing semantic entropy. Merge-aware systems retain the insight of combination while bounding it with persistence and provenance.

The comparison with remix culture sharpened the critique further. Remix excels at production and propagation but fails as a memory system. Feed platforms implicitly adopt remix logic as a default, treating all artifacts as disposable and endlessly recontextualizable. Gallery-first systems reject this as a foundation. They allow remix, but only as an explicit event referencing a stable object, preserving lineage rather than erasing it.

The design corollaries made these abstractions visible. The persistent presence of prompts such as “Write a note,” the mandatory display of like and comment buttons, and the forced insertion of algorithmic content are not superficial annoyances. They are visible proofs that the system does not recognize a completed, silent, curated state as legitimate. By contrast, the calm coherence of archived views demonstrates that an alternative projection is already possible, but philosophically disallowed.

The formal transition rules completed the argument by showing that coherence can be enforced by construction. By explicitly separating gallery, interaction, and annotation states, and by requiring deliberate transitions between them, the system encodes restraint as a property of the architecture rather than a demand placed on users. Accidental engagement becomes impossible. Meaning-bearing actions become intentional. Silence becomes stable.

Taken together, these results support a strong conclusion. Meaning cannot survive in systems that treat motion as value and reaction as perception. Memory cannot form where every surface demands output. Identity cannot persist where objects lack stable reference.

A gallery-first, merge-aware system does not promise novelty, virality, or growth. It promises something more fundamental: the possibility that artifacts can settle, histories can accrete, and identity can remain coherent over time. This is not a nostalgic appeal to slower media, but a structural correction to architectures that confuse optimization with understanding.

The choice, ultimately, is architectural. Systems can be built to extract engagement, or they can be built to conserve meaning. They cannot do both simultaneously. The frustration that motivated this paper is therefore not a personal complaint. It is an early warning signal of a deeper incompatibility between human sense-making and engagement-driven infrastructure.

Meaning cannot be optimized into existence. It must be allowed to remain.

\newpage
\appendix
\section{Formal Appendix: Minimal Conditions for Semantic Drift}

This appendix formalizes the claim that feed-based architectures necessarily exhibit unbounded semantic drift, while gallery-first, merge-aware systems necessarily permit semantic accumulation. The argument is constructive and relies only on minimal structural assumptions.

\subsection{Preliminaries}

Let $\mathcal{C}$ denote the set of content instances presented to a user over time. Let $\mathcal{E}$ denote the set of events, and let $\mathcal{R}$ denote the set of reactions.

We distinguish between two architectural regimes.

In a \emph{feed-based system}, content instances are primary, reactions bind to instances, and ordering is chronological.

In a \emph{merge-aware system}, semantic objects are primary, events reference objects, and feeds are projections.

\subsection{Assumptions for Feed-Based Systems}

A feed-based system satisfies the following conditions.

First, each content instance $c \in \mathcal{C}$ has a unique identity even when two instances have identical payloads.

Second, reactions bind to instances rather than to payload equivalence classes. That is, if $c_1 \neq c_2$, then reactions to $c_1$ and $c_2$ are disjoint, regardless of content similarity.

Third, the system admits repeated introduction of semantically equivalent content instances over time.

Fourth, there exists no canonical merge operator $\mu$ such that semantically equivalent instances are unified into a single referent.

These assumptions are satisfied by all mainstream engagement-driven feed architectures.

\subsection{Definition of Semantic Drift}

Let $M(c)$ denote the meaning associated with content instance $c$, defined as the aggregate of reactions, annotations, and contextual associations accumulated over time.

Semantic drift occurs if, for a fixed payload $p$, the variance of $M(c)$ over instances $c$ carrying $p$ grows unbounded as the number of instances increases.

Formally, drift occurs if
\[
\lim_{n \to \infty} \mathrm{Var}\left( \{ M(c_i) \mid c_i \text{ carries } p \} \right) = \infty.
\]

\subsection{Drift Theorem for Feed-Based Systems}

\textbf{Theorem 1.}  
In any feed-based system satisfying the assumptions above, semantic drift is inevitable for any payload that is reintroduced infinitely often.

\emph{Proof.}  
Let $\{c_1, c_2, \dots\}$ be an infinite sequence of content instances carrying identical payload $p$.

By assumption, reactions bind to instances, not to payloads. Therefore, each $c_i$ accumulates a distinct reaction history $M(c_i)$.

Because instances appear in different temporal, social, and contextual neighborhoods, the distributions of reactions differ across instances. No mechanism exists to reconcile or unify these histories.

Since the system admits repeated introduction without merge, the set $\{M(c_i)\}$ grows without bound in cardinality and diversity. Therefore, the variance of meaning associated with payload $p$ increases monotonically with the number of instances.

Hence semantic drift is inevitable. \hfill $\square$

\subsection{Assumptions for Merge-Aware Systems}

A merge-aware system satisfies the following conditions.

First, there exists a set of semantic objects $\mathcal{O}$ such that all events reference objects rather than instances.

Second, there exists an equivalence relation $\sim$ over payloads such that semantically identical payloads reference the same object.

Third, reactions bind to objects via events.

Fourth, there exists a merge operator $\mu$ that is associative, commutative, and idempotent over events referencing the same object.

\subsection{Definition of Semantic Accumulation}

Let $M(o)$ denote the meaning associated with object $o$, defined as the aggregation of all reactions and annotations across events referencing $o$.

Semantic accumulation occurs if repeated references to $o$ increase $|M(o)|$ without increasing the number of distinct semantic referents.

\subsection{Accumulation Theorem for Merge-Aware Systems}

\textbf{Theorem 2.}  
In any merge-aware system satisfying the assumptions above, repeated reference to a semantic object results in accumulation rather than drift.

\emph{Proof.}  
Let $o \in \mathcal{O}$ be a semantic object with payload $p$. Let $\{e_1, e_2, \dots\}$ be an infinite sequence of events referencing $o$.

By construction, all reactions generated by these events bind to $o$. The merge operator $\mu$ ensures that event histories are consolidated without duplication or loss.

Therefore, the meaning function $M(o)$ grows monotonically while the referent set remains singleton. No variance across semantic identities is introduced.

Hence repetition strengthens meaning rather than fragmenting it. \hfill $\square$

\subsection{Corollary: Incompatibility of Feed Optimization and Semantic Conservation}

\textbf{Corollary.}  
Any system that optimizes for engagement via instance-level metrics while lacking a merge operator over semantic equivalence classes cannot conserve meaning over time.

\emph{Justification.}  
Instance-level optimization incentivizes duplication, resurfacing, and recontextualization without consolidation. Without merge, these operations necessarily increase semantic variance.

Therefore, semantic conservation and engagement optimization are structurally incompatible unless the latter is subordinated to object-level identity.

\subsection{Interpretation}

These results are intentionally minimal. They do not depend on user behavior, moderation quality, or recommendation accuracy. They follow directly from structural choices about identity, binding, and merge.

Feed-based drift is not an accident. It is a theorem.

Merge-aware accumulation is not a feature. It is a consequence.

This completes the formal argument that coherence must be designed, not tuned.

\newpage
\begin{thebibliography}{99}

\bibitem{FauconnierTurner2002}
G. Fauconnier and M. Turner.
\newblock \emph{The Way We Think: Conceptual Blending and the Mind's Hidden Complexities}.
\newblock Basic Books, 2002.

\bibitem{Turner2014}
M. Turner.
\newblock \emph{The Origin of Ideas: Blending, Creativity, and the Human Spark}.
\newblock Oxford University Press, 2014.

\bibitem{Brand2018}
S. Brand.
\newblock \emph{How Buildings Learn: What Happens After They’re Built}.
\newblock Penguin Books, 2018.
\newblock Originally published 1994.

\bibitem{Lamport1998}
L. Lamport.
\newblock \emph{The Part-Time Parliament}.
\newblock ACM Transactions on Computer Systems, 16(2), 1998.

\bibitem{Fowler2017}
M. Fowler.
\newblock \emph{Event Sourcing}.
\newblock martinfowler.com, 2017.
\newblock Online essay.

\bibitem{Packer2013}
J. Packer and S. B. Crofts Wiley (eds.).
\newblock \emph{Communication Matters: Materialist Approaches to Media, Mobility and Networks}.
\newblock Routledge, 2013.

\bibitem{Manovich2001}
L. Manovich.
\newblock \emph{The Language of New Media}.
\newblock MIT Press, 2001.

\bibitem{Lessig2008}
L. Lessig.
\newblock \emph{Remix: Making Art and Commerce Thrive in the Hybrid Economy}.
\newblock Penguin Press, 2008.

\bibitem{Benkler2006}
Y. Benkler.
\newblock \emph{The Wealth of Networks}.
\newblock Yale University Press, 2006.

\bibitem{Tufte2001}
E. R. Tufte.
\newblock \emph{The Visual Display of Quantitative Information}.
\newblock Graphics Press, 2001.

\bibitem{Norman2013}
D. A. Norman.
\newblock \emph{The Design of Everyday Things}.
\newblock Basic Books, Revised Edition, 2013.

\bibitem{Rogers2013}
R. Rogers.
\newblock \emph{Digital Methods}.
\newblock MIT Press, 2013.

\bibitem{Zuboff2019}
S. Zuboff.
\newblock \emph{The Age of Surveillance Capitalism}.
\newblock PublicAffairs, 2019.

\bibitem{OReilly2017}
T. O’Reilly.
\newblock \emph{WTF?: What’s the Future and Why It’s Up to Us}.
\newblock Harper Business, 2017.

\bibitem{Simon1962}
H. A. Simon.
\newblock \emph{The Architecture of Complexity}.
\newblock Proceedings of the American Philosophical Society, 106(6), 1962.

\bibitem{Floridi2014}
L. Floridi.
\newblock \emph{The Fourth Revolution: How the Infosphere Is Reshaping Human Reality}.
\newblock Oxford University Press, 2014.

\bibitem{CSCW2019}
M. Muller et al.
\newblock \emph{Designing for Meaningful Social Interaction}.
\newblock Proceedings of CSCW, 2019.

\bibitem{Hutchins1995}
E. Hutchins.
\newblock \emph{Cognition in the Wild}.
\newblock MIT Press, 1995.

\bibitem{HofstadterSander2013}
D. R. Hofstadter and E. Sander.
\newblock \emph{Surfaces and Essences: Analogy as the Fuel and Fire of Thinking}.
\newblock Basic Books, 2013.

\end{thebibliography}

\end{document}
